\newpage
\section{Troisième génération : DHT structurée inspirée de Chord}

Cette troisième génération correspond au passage vers une DHT réellement structurée. L'idée n'est plus de compenser les limites d'un réseau non structuré par du flooding ou par une connaissance partielle des pairs, mais de reposer sur un espace d'identifiants commun et sur un routage déterministe.

Plusieurs protocoles illustrent cette famille de DHT structurées. \textit{Chord} organise les n\oe uds sur un anneau logique et utilise des sauts logarithmiques à l'aide d'une \textit{finger table} pour retrouver le responsable d'une clé. \textit{Pastry} repose sur un routage par préfixes : à chaque étape, le message progresse vers un n\oe ud dont l'identifiant partage un préfixe plus long avec la clé recherchée. \textit{Kademlia}, de son côté, structure la recherche autour d'une distance XOR entre identifiants et choisit à chaque saut des n\oe uds de plus en plus proches de la cible selon cette métrique. D'autres approches comme \textit{CAN} répartissent enfin l'espace des clés dans un espace cartésien virtuel où le routage se fait de zone en zone entre voisins.

Dans la suite du rapport, nous nous concentrons toutefois sur \textit{Chord}, qui sert de base à l'implémentation présentée. Les autres protocoles sont évoqués ici comme points de comparaison pour situer notre travail parmi les principales DHT structurées.

\subsection{Vue d'ensemble de l'architecture}

Cette troisi\`eme impl\'ementation adopte une architecture DHT structur\'ee inspir\'ee de Chord. Contrairement aux g\'en\'erations pr\'ec\'edentes, il n'existe ni \emph{coordinator} central ni \emph{flooding} : chaque n\oe ud s'auto-organise dans un anneau logique et route de mani\`ere d\'eterministe les requ\^etes vers le n\oe ud responsable, sans connaissance globale du r\'eseau.

Cette troisi\`eme impl\'ementation repose sur deux concepts fondamentaux :
\begin{itemize}
    \item Un \textbf{anneau logique} : chaque n\oe ud occupe une position fixe sur un anneau de taille \(2^6 = 64\), calcul\'ee par hachage de son adresse. Les cl\'es sont distribu\'ees sur cet anneau via \emph{consistent hashing}.
    \item Une \textbf{finger table} : table de routage locale de \(M = 6\) entr\'ees permettant \`a chaque n\oe ud de router les requ\^etes en \(O(\log N)\) messages sans connaissance globale du r\'eseau.
\end{itemize}
La coh\'erence de l'anneau est assur\'ee par un service de stabilisation p\'eriodique qui corrige les pointeurs \emph{successor} et \emph{predecessor} et maintient une liste de successeurs de secours pour la tol\'erance aux pannes.

\subsubsection{Principe de fonctionnement}

Chaque n\oe ud et chaque cl\'e re\c coivent un identifiant sur \(m\) bits (souvent via SHA-1), plac\'e sur un anneau modulo \(2^m\). Une cl\'e \(k\) est stock\'ee sur son \emph{successor} : le premier n\oe ud rencontr\'e dans le sens horaire dont l'identifiant est sup\'erieur ou \`a \(k\).

Cette r\`egle donne trois propri\'et\'es utiles :
\begin{itemize}
\item \textbf{R\'epartition} : le \emph{consistent hashing} distribue les cl\'es de mani\`ere quasi uniforme.
\item \textbf{Scalabilit\'e} : l'ajout/retrait d'un n\oe ud ne d\'eplace qu'une fraction des cl\'es contrairement au \texttt{hash(key) \% N} des g\'en\'erations pr\'ec\'edentes.
\item \textbf{D\'ecentralisation} : aucun serveur central n'est requis.
\end{itemize}

\subsubsection{Routage efficace : Finger Table}

Pour \`eviter un parcours lin\'eaire de l'anneau, chaque n\oe ud maintient une \textbf{finger table} de \(m\) entr\'ees. L'entr\'ee \(i\) pointe vers le premier n\oe ud joignable \`a partir de \(n + 2^{i-1}\).

\[
\text{finger}[i].\text{start} = (n + 2^{i-1}) \bmod 2^m
\]

Cette structure permet des sauts exponentiels et un co\^ut moyen de recherche en \(O(\log N)\) messages : le routage consiste \`a s\'electionner le finger le plus proche pr\'ec\'edant la cl\'e cible (\texttt{closestPrecedingFinger}), et \`a d\'el\'eguer la requ\^ete \`a ce n\oe ud, jusqu'\`a atteindre le \emph{successor} de la cl\'e.


\subsubsection{Op\'eration PUT}

Lorsqu'un n\oe ud re\c{c}oit une requ\^ete \texttt{PUT} pour une paire \((key, value)\), il calcule le \emph{hash} de la cl\'e et appelle \emph{lookup} pour identifier le \emph{successor} responsable. Si ce \emph{successor} est le n\oe ud lui-m\^eme, la paire est stock\'ee localement. Sinon, la requ\^ete est rout\'ee vers le n\oe ud d\'esign\'e via la \emph{finger table}. Le message transite ainsi de n\oe ud en n\oe ud jusqu'\`a atteindre le responsable, qui effectue le stockage final.

\textbf{Nombre total de messages \'echang\'es :} \(O(\log N)\)

\subsubsection{Op\'eration GET}

La lecture suit exactement la m\^eme logique que le \texttt{PUT}. Le n\oe ud recevant la requ\^ete \texttt{GET} calcule le \emph{hash} de la cl\'e et identifie le \emph{successor} responsable via \emph{lookup}. Si c'est lui-m\^eme, il retourne la valeur depuis son \emph{store} local, ou \texttt{NOT\_FOUND} si la cl\'e est absente. Sinon, la requ\^ete est rout\'ee vers le n\oe ud responsable, qui r\'epond directement \`a l'\'emetteur original.

\begin{figure}[!ht]
    \centering
    \includegraphics[width=0.6\textwidth,height=0.45\textheight,keepaspectratio]{prism-uploads/GetCHORD.drawio.pdf}
    \caption{D\'eroulement d'une op\'eration \texttt{GET} dans l'architecture Chord.}
    \label{fig:chord_get}
\end{figure}

\textbf{Nombre total de messages \'echang\'es :} \(O(\log N)\)


\subsubsection{Int\'egration d'un nouveau n\oe ud (JOIN)}

Lorsqu'un nouveau n\oe ud d\'emarre, il contacte un n\oe ud existant pass\'e en param\`etre (\texttt{bootstrapPeer}) et appelle \texttt{remoteCallFindSuccessor} pour identifier sa position sur l'anneau. Une fois son \emph{successor} d\'etermin\'e, il lui envoie un message \texttt{REQUEST\_KEYS} pour r\'ecup\'erer les cl\'es dont il devient responsable. Contrairement \`a la premi\`ere g\'en\'eration, ce transfert est cibl\'e : seules les cl\'es appartenant \`a la plage \((\text{pr\'ed\'ecesseur du successor}, \text{id du nouveau n\oe ud}]\) sont d\'eplac\'ees. Le service de stabilisation se charge ensuite de propager progressivement la connaissance du nouveau n\oe ud aux autres membres de l'anneau via \texttt{notify}.

\subsubsection{Retrait d'un n\oe ud (LEAVE)}

Lors d'un d\'epart volontaire, le \texttt{ShutdownHook} d\'eclenche la proc\'edure de sortie propre. Le n\oe ud transf\`ere d'abord l'int\'egralit\'e de son \emph{store} local \`a son \emph{successor} via \texttt{transferAllKeysTo}. Il envoie ensuite un message \texttt{LEAVE\_NOTIFY} \`a son \emph{successor} et \`a son \emph{predecessor} pour qu'ils mettent \`a jour leurs pointeurs et maintiennent la coh\'erence de l'anneau.

\subsubsection{D\'etection des pannes et stabilisation}

Le service de stabilisation (\texttt{StabilizeService}) s'ex\'ecute toutes les 2 secondes. \`A chaque cycle, il v\'erifie la vivacit\'e du \emph{successor} via \texttt{remoteCallGetPredecessor} : si celui-ci ne r\'epond pas, il cherche un rempla\c{c}ant dans la liste de successeurs de secours, puis dans la \emph{finger table}. Le \emph{predecessor} est v\'erifi\'e tous les 3 cycles. En cas de panne, il est remis \`a \texttt{null} et l'anneau se reconfigure progressivement. La \emph{finger table} est quant \`a elle mise \`a jour entr\'ee par entr\'ee \`a chaque cycle via \texttt{fixFingers}.

\subsubsection{Avantages et limites de l'approche Chord}

L'architecture Chord pr\'esente plusieurs avantages majeurs par rapport aux g\'en\'erations pr\'ec\'edentes. Le \emph{consistent hashing} garantit une distribution quasi uniforme des cl\'es et limite les transferts lors des changements de topologie : un \emph{JOIN} ou \emph{LEAVE} ne d\'eplace qu'une fraction des cl\'es, contre une redistribution massive dans la premi\`ere g\'en\'eration. Le routage d\'eterministe via la \emph{finger table} offre un \emph{lookup} en \(O(\log N)\) messages, sans \emph{flooding} ni \emph{coordinator} central. Enfin, la liste de successeurs de secours et la stabilisation p\'eriodique assurent une tol\'erance aux pannes que les g\'en\'erations pr\'ec\'edentes ne permettaient pas.

Cependant, cette architecture pr\'esente \'egalement des limites. La stabilisation p\'eriodique introduit un d\'elai de convergence : dans une p\'eriode de \emph{churn} intense, l'anneau peut temporairement \^etre dans un \'etat incoh\'erent, rendant certaines cl\'es introuvables le temps que les pointeurs se corrigent. Par ailleurs, notre impl\'ementation utilise un espace d'adressage de taille \(2^6 = 64\), ce qui limite le nombre de n\oe uds distincts pouvant coexister sur l'anneau. Enfin, comme dans la premi\`ere g\'en\'eration, l'absence de r\'eplication implique qu'une panne non d\'etect\'ee avant la stabilisation peut entra\^iner une perte temporaire de donn\'ees.

\input{sections/03b_chord_implementation}